\documentclass[10pt,a4paper]{amsart}
\usepackage{amsmath,amssymb,amsthm,graphicx}
\usepackage{verbatim}
\newenvironment{reference}{\comment}{\endcomment}
%\newenvironment{reference}{}{}
\newenvironment{slogan}{\comment}{\endcomment}
\newenvironment{history}{\comment}{\endcomment}
% For commutative diagrams we use Xy-pic
\usepackage[all]{xy}

% We use 2cell for 2-commutative diagrams.
\xyoption{2cell}
% Theorem environments.
%
\theoremstyle{plain}
\newtheorem{theorem}[subsection]{Theorem}
\newtheorem{proposition}[subsection]{Proposition}
\newtheorem{lemma}[subsection]{Lemma}

\theoremstyle{definition}
\newtheorem{definition}[subsection]{Definition}
\newtheorem{example}[subsection]{Example}
\newtheorem{exercise}[subsection]{Exercise}
\newtheorem{situation}[subsection]{Situation}

\theoremstyle{remark}
\newtheorem{remark}[subsection]{Remark}
\newtheorem{remarks}[subsection]{Remarks}

\numberwithin{equation}{subsection}

% Macros
%
\def\lim{\mathop{\mathrm{lim}}\nolimits}
\def\colim{\mathop{\mathrm{colim}}\nolimits}
\def\Spec{\mathop{\mathrm{Spec}}}
\def\Hom{\mathop{\mathrm{Hom}}\nolimits}
\def\Ext{\mathop{\mathrm{Ext}}\nolimits}
\def\SheafHom{\mathop{\mathcal{H}\!\mathit{om}}\nolimits}
\def\SheafExt{\mathop{\mathcal{E}\!\mathit{xt}}\nolimits}
\def\Sch{\mathit{Sch}}
\def\Mor{\mathop{\mathrm{Mor}}\nolimits}
\def\Ob{\mathop{\mathrm{Ob}}\nolimits}
\def\Sh{\mathop{\mathit{Sh}}\nolimits}
\def\NL{\mathop{N\!L}\nolimits}
\def\CH{\mathop{\mathrm{CH}}\nolimits}
\def\proetale{{pro\text{-}\acute{e}tale}}
\def\etale{{\acute{e}tale}}
\def\QCoh{\mathit{QCoh}}
\def\Ker{\mathop{\mathrm{Ker}}}
\def\Im{\mathop{\mathrm{Im}}}
\def\Coker{\mathop{\mathrm{Coker}}}
\def\Coim{\mathop{\mathrm{Coim}}}

% Boxtimes
%
\DeclareMathSymbol{\boxtimes}{\mathbin}{AMSa}{"02}

%
% Macros for moduli stacks/spaces
%
\def\QCohstack{\mathcal{QC}\!\mathit{oh}}
\def\Cohstack{\mathcal{C}\!\mathit{oh}}
\def\Spacesstack{\mathcal{S}\!\mathit{paces}}
\def\Quotfunctor{\mathrm{Quot}}
\def\Hilbfunctor{\mathrm{Hilb}}
\def\Curvesstack{\mathcal{C}\!\mathit{urves}}
\def\Polarizedstack{\mathcal{P}\!\mathit{olarized}}
\def\Complexesstack{\mathcal{C}\!\mathit{omplexes}}
% \Pic is the operator that assigns to X its picard group, usage \Pic(X)
% \Picardstack_{X/B} denotes the Picard stack of X over B
% \Picardfunctor_{X/B} denotes the Picard functor of X over B
\def\Pic{\mathop{\mathrm{Pic}}\nolimits}
\def\Picardstack{\mathcal{P}\!\mathit{ic}}
\def\Picardfunctor{\mathrm{Pic}}
\def\Deformationcategory{\mathcal{D}\!\mathit{ef}}
\usepackage{cleveref}
\crefname{originallemma}{original lemma}{original lemmas}
\Crefname{originallemma}{Original lemma}{Original lemmas}
\makeatletter
\newlabel{lemma-DM}{{lemma-DM}{source}{}{sourcepointer-1}{}}
\newlabel{lemma-DM@cref}{{[originallemma][0][]lemma-DM}{[1][0][]source}{}{}{}}
\makeatother
\crefname{originallemma}{original lemma}{original lemmas}
\Crefname{originallemma}{Original lemma}{Original lemmas}
\makeatletter
\newlabel{lemma-compose-after-separated}{{lemma-compose-after-separated}{source}{}{sourcepointer-2}{}}
\newlabel{lemma-compose-after-separated@cref}{{[originallemma][0][]lemma-compose-after-separated}{[1][0][]source}{}{}{}}
\makeatother
\crefname{originalreference}{original reference}{original references}
\Crefname{originalreference}{Original reference}{Original references}
\makeatletter
\newlabel{spaces-morphisms-lemma-etale-permanence}{{spaces-morphisms-lemma-etale-permanence}{source}{}{sourcepointer-3}{}}
\newlabel{spaces-morphisms-lemma-etale-permanence@cref}{{[originalreference][0][]spaces-morphisms-lemma-etale-permanence}{[1][0][]source}{}{}{}}
\makeatother
\title{A morphism of algebraic stacks etale over a common base is etale}
\author{The Stacks Project Authors}
\date{stacks-project:stacks-morphisms.tex:local-snapshot-20261012; GNU FDL 1.2 or later}
\begin{document}
\maketitle
\noindent\textit{Editorial note.} All delivered mathematical text is the author's own, copied from the pinned Stacks Project chapter source sources/stacks-algebraic-stacks-morphisms/pinned-original-source.tex; no statement or proof step was rewritten, reordered or removed.\par\medskip
\noindent\textit{Editorial note.} Preamble provenance: the Stacks Project keeps one single shared preamble for every chapter; this card takes it directly from the staged sources/stacks-algebraic-stacks-morphisms/preamble.tex as three verbatim spans (lines 8-12, 14-18 and 154-217), so no derived file was written and nothing was added or redefined. The wrapper class and the base packages are supplied by the assembler.\par\medskip
\noindent\textit{Editorial note.} Bibliography provenance: neither the counted statement nor its proof contains a citation (this lane rescans the whole span before writing the card), so the card carries no bibliography and no entry was reconstructed.\par\medskip
\noindent\textit{Editorial note.} Selection provenance: this card was opened under the file-based control contract (CONTROL.md), with the mechanical pre-card gate run on its statement and proof interval before the card was written.\par\medskip
\noindent\textit{Allowed standard prerequisites (editorial).}\begin{itemize}
\item The statement that a morphism of algebraic stacks between separated or DM stacks over a base is DM. Invoked at the start of the counted proof (source lines 8205-8206). Makes the morphism DM so that it admits etale presentations by algebraic spaces.
\item The existence of etale presentations for DM morphisms of algebraic stacks. Invoked twice inside the counted proof (source lines 8209-8215). Produces the algebraic spaces V and U presenting the two stacks over the base and over the target.
\item The version for algebraic spaces of the etale permanence statement. Invoked at the end of the counted proof (source lines 8225-8228) as a cross-chapter pointer. Gives etaleness of the induced morphism of presenting spaces.
\end{itemize}\medskip
\ifdefined\hypertarget\hypertarget{sourcepointer-1}{}\fi
\noindent\textit{Source reference.} Original lemma lemma-DM: sources/stacks-algebraic-stacks-morphisms/pinned-original-source.tex, source lines 7800--7800. Navigation pointer only: the counted text refers to the Stacks item lemma-DM, which belongs to another chapter of the project and is not part of this card. The pointer carries the target's own label as its number because this chapter's staged snapshot does not contain the other chapter's tag table; its content is not reproduced here.\par\medskip
\ifdefined\hypertarget\hypertarget{sourcepointer-2}{}\fi
\noindent\textit{Source reference.} Original lemma lemma-compose-after-separated: sources/stacks-algebraic-stacks-morphisms/pinned-original-source.tex, source lines 1089--1089. Navigation pointer only: the counted text refers to the Stacks item lemma-compose-after-separated, which belongs to another chapter of the project and is not part of this card. The pointer carries the target's own label as its number because this chapter's staged snapshot does not contain the other chapter's tag table; its content is not reproduced here.\par\medskip
\ifdefined\hypertarget\hypertarget{sourcepointer-3}{}\fi
\noindent\textit{Source reference.} Original reference spaces-morphisms-lemma-etale-permanence: sources/stacks-algebraic-stacks-morphisms/pinned-original-source.tex, source lines 4653--4653. Navigation pointer only: the counted text refers to the Stacks item spaces-morphisms-lemma-etale-permanence, which belongs to another chapter of the project and is not part of this card. The pointer carries the target's own label as its number because this chapter's staged snapshot does not contain the other chapter's tag table; its content is not reproduced here.\par\medskip
\begin{lemma}
\label{lemma-etale-permanence}
Let $\mathcal{X}, \mathcal{Y}$ be algebraic stacks \'etale over
an algebraic stack $\mathcal{Z}$. Any morphism
$\mathcal{X} \to \mathcal{Y}$ over $\mathcal{Z}$ is \'etale.
\end{lemma}

\begin{proof}
The morphism $\mathcal{X} \to \mathcal{Y}$ is DM by
Lemma \ref{lemma-compose-after-separated}.
Let $W \to \mathcal{Z}$ be a surjective smooth morphism
whose source is an algebraic space. Let
$V \to \mathcal{Y} \times_\mathcal{Z} W$ be a surjective
\'etale morphism whose source is an algebraic space
(Lemma \ref{lemma-DM}).
Let $U \to \mathcal{X} \times_\mathcal{Y} V$ be a surjective
\'etale morphism whose source is an algebraic space
(Lemma \ref{lemma-DM}).
Then
$$
U \longrightarrow \mathcal{X} \times_\mathcal{Z} W
$$
is surjective \'etale as the composition of
$U \to \mathcal{X} \times_\mathcal{Y} V$
and the base change of $V \to \mathcal{Y} \times_\mathcal{Z} W$
by $\mathcal{X} \times_\mathcal{Z} W \to \mathcal{Y} \times_\mathcal{Z} W$.
Hence it suffices to show that $U \to W$ is \'etale.
Since $U \to W$ and $V \to W$ are \'etale because
$\mathcal{X} \to \mathcal{Z}$ and $\mathcal{Y} \to \mathcal{Z}$
are \'etale, this follows from the version of the lemma for
algebraic spaces, namely
Morphisms of Spaces, Lemma \ref{spaces-morphisms-lemma-etale-permanence}.
\end{proof}

\end{document}
